\chapter{Priapism}

\section*{Definition}
Priapism is a persistent, often painful erection lasting longer than 4 hours unrelated to or persisting beyond sexual stimulation. Ischemic (low-flow, veno-occlusive) priapism accounts for \textgreater{}95\% of cases; non-ischemic (high-flow, arterial) priapism is less common and typically painless.

\section*{What Happens in Your Body}
In ischemic priapism, venous outflow obstruction or failure of detumescence mechanisms leads to a compartment syndrome of the corpora cavernosa: hypoxia, acidosis, and eventual smooth-muscle necrosis if untreated. Non-ischemic priapism results from uncontrolled arterial inflow due to cavernosal artery laceration (usually after perineal trauma) bypassing normal regulatory mechanisms.

\section*{Causes}
\begin{itemize}
  \item \textbf{Hematologic:} Sickle cell disease (most common cause in children), thalassemia, polycythemia vera, essential thrombocythemia, multiple myeloma
  \item \textbf{Pharmacologic:} Intracavernosal injections (alprostadil, papaverine, phentolamine), trazodone, antipsychotics (especially risperidone, olanzapine), sildenafil/tadalafil (rare), antihypertensives (prazosin, hydralazine), anticoagulants, cocaine, cannabis
  \item \textbf{Neoplastic:} Penile, prostatic, urethral, or bladder malignancies causing venous infiltration; metastatic disease
  \item \textbf{Neurologic:} Spinal cord injury, cauda equina syndrome, general anesthesia
  \item \textbf{Trauma:} Perineal or penile trauma leading to cavernosal artery rupture (high-flow priapism)
  \item \textbf{Idiopathic:} Up to 30\% of ischemic cases have no identifiable cause
\end{itemize}

\section*{When to See a Doctor}
See your doctor if:
\begin{itemize}
  \item Erection persists for \textgreater{}4 hours, regardless of pain
  \item Pain in the penile shaft or perineum with rigid corpora cavernosa and a soft glans
  \item Known sickle cell disease with prolonged erection
  \item History of perineal trauma with painless, persistent erection
\end{itemize}

\section*{Related Symptoms}
\begin{itemize}
  \item Penile pain, lower urinary tract symptoms, perineal fullness, dysuria, erectile dysfunction (as a sequela of delayed treatment), priapism recurrence (stuttering priapism)
\end{itemize}
\textbf{Important:} Priapism is a urologic emergency; the risk of irreversible erectile dysfunction increases significantly after 12--24 hours of ischemia, with fibrosis of the corpus cavernosum becoming histologically evident beyond 48 hours.

\section*{Self-Care / Home Management}
\begin{itemize}
  \item Emergent urologic evaluation is required---no safe home management exists for ischemic priapism
  \item Apply ice packs to the perineum and thighs during transport to the emergency department
  \item For stuttering (recurrent) priapism in sickle cell disease, encourage hydration and oral analgesics while seeking care
  \item Empty bladder if possible (voiding may provide minor relief but should not delay treatment)
\end{itemize}

\section*{How Your Doctor Will Evaluate You}
\begin{itemize}
  \item Corporal blood gas sampling: Ischemic priapism shows pO\(\_2\) \textless{}30 mmHg, pCO\(\_2\) \textgreater{}60 mmHg, pH \textless{}7.25; non-ischemic shows arterial blood gas values
  \item Penile Doppler ultrasound: Assess cavernosal arterial flow---absent or minimal flow in ischemic priapism; high-velocity flow with turbulent pattern in non-ischemic priapism
  \item Laboratory studies: CBC with sickle cell screen, coagulation profile, reticulocyte count, urine toxicology, hemoglobin electrophoresis if indicated
  \item MRI of penis if concern for neoplastic infiltration or trauma-related pseudoaneurysm
\end{itemize}

\section*{Treatment Options}
\begin{itemize}
  \item \textbf{Ischemic (low-flow):} Aspiration and irrigation of corpora cavernosa with saline; intracavernosal injection of phenylephrine (100--500 \(\mu\)g, repeated every 3--5 minutes); surgical shunting (Winter, Al-Ghorab, or T--shunt) if medical management fails; exchange transfusion for sickle cell disease
  \item \textbf{Non-ischemic (high-flow):} Conservative management with observation or perineal ice packs; selective arterial embolization (gelatin foam or coils) if persistent; surgical ligation if embolization fails
  \item \textbf{Stuttering priapism:} Preventive therapy with gonadotropin-releasing hormone agonists, oral ketoconazole, baclofen, or PDE-5 inhibitors (low-dose daily tadalafil) to suppress recurrent episodes
  \item \textbf{Follow-up:} Urology referral; counseling on erectile dysfunction risk; psychological support for recurrent episodes
\end{itemize}

\section*{References}
\begin{thebibliography}{99}
\bibitem{priapism:priapism1} Montague DK, Jarow J, Broderick GA, et al. "American Urological Association guideline on the management of priapism." J Urol. 2003;170(4 Pt 1):1318-1324.
\bibitem{priapism:priapism2} Burnett AL. "Priapism: current principles and practice." Urol Clin North Am. 2007;34(4):631-642.
\bibitem{priapism:priapism3} Broderick GA, Kadioglu A, Bivalacqua TJ, et al. "Priapism: pathogenesis, epidemiology, and management." J Sex Med. 2010;7(1 Pt 2):476-500.
\bibitem{priapism:priapism4} Bivalacqua TJ, Hellstrom WJ. "Priapism." In: Campbell-Walsh Urology. 12th ed. Elsevier; 2021.
\bibitem{priapism:priapism5} Kovac JR, Mak SK, Garcia MM, et al. "A pathophysiology-based approach to the management of early priapism." Asian J Androl. 2013;15(1):20-26.
\bibitem{priapism:priapism6} Burnett AL, Bivalacqua TJ. "Priapism: new concepts in medical and surgical management." Urol Clin North Am. 2021;48(4):529-540.
\end{thebibliography}
